% ═══════════════════════════════════════════════════════════════
% MT-03c-ML: Minitest Los colores - MUSTERLÖSUNG
% Spanisch Klasse 7 | Sequenz 3: Mi casa
% ═══════════════════════════════════════════════════════════════

\documentclass[11pt, a4paper]{article}

\usepackage[utf8]{inputenc}
\usepackage[T1]{fontenc}
\usepackage[ngerman]{babel}
\usepackage{helvet}
\renewcommand{\familydefault}{\sfdefault}

\usepackage[margin=2cm]{geometry}
\usepackage{enumitem}
\usepackage{tabularx}
\usepackage{booktabs}

\usepackage{xcolor}
\usepackage{tcolorbox}
\tcbuselibrary{skins}

\usepackage{fancyhdr}
\usepackage{lastpage}
\usepackage{needspace}

% FARBEN
\definecolor{spanisch-color}{HTML}{c41e3a}
\definecolor{grau-color}{HTML}{888888}
\definecolor{hellgrau-color}{HTML}{f5f5f5}
\definecolor{loesung-color}{HTML}{c62828}

% Farbkästchen
\definecolor{rojo}{HTML}{e53935}
\definecolor{azul}{HTML}{1e88e5}
\definecolor{verde}{HTML}{43a047}
\definecolor{amarillo}{HTML}{fdd835}
\definecolor{negro}{HTML}{212121}
\definecolor{blanco}{HTML}{ffffff}
\definecolor{marron}{HTML}{795548}
\definecolor{gris}{HTML}{9e9e9e}
\definecolor{rosa}{HTML}{ec407a}
\definecolor{naranja}{HTML}{fb8c00}

\colorlet{spanisch}{spanisch-color}
\colorlet{grau}{grau-color}
\colorlet{hellgrau}{hellgrau-color}
\colorlet{loesung}{loesung-color}

\newcommand{\fachfarbe}{spanisch}

\newcommand{\thema}{Minitest: Los colores}
\newcommand{\fach}{Spanisch}
\newcommand{\klasse}{7}
\newcommand{\maxpunkte}{30}
\newcommand{\zeit}{15 Minuten}
\newcommand{\datum}{\today}

\pagestyle{fancy}
\fancyhf{}
\renewcommand{\headrulewidth}{0.4pt}
\renewcommand{\footrulewidth}{0.4pt}

\lhead{\fach{} -- Klasse \klasse}
\chead{\textcolor{loesung}{\textbf{MT-03c MUSTERLÖSUNG}}}
\rhead{\datum}

\lfoot{\zeit}
\cfoot{}
\rfoot{Seite \thepage\ von \pageref{LastPage}}

\newcommand{\punktebox}[1]{\hfill\fbox{\small \_/#1}}
\newcommand{\luecke}[1]{\rule{#1}{0.4pt}}
\newcommand{\lsg}[1]{\textcolor{loesung}{\textbf{#1}}}
\newcommand{\farbbox}[1]{\colorbox{#1}{\hspace{0.8em}\rule{0pt}{0.8em}}}

\newtcolorbox{infobox}{
    colback=hellgrau,
    colframe=\fachfarbe,
    boxrule=1pt,
    arc=0pt,
    left=8pt, right=8pt, top=8pt, bottom=8pt
}

\newtcolorbox{regelbox}{
    colback=white,
    colframe=\fachfarbe,
    boxrule=1pt,
    arc=0pt,
    left=8pt, right=8pt, top=8pt, bottom=8pt
}

\newenvironment{aufgabe}[1]{%
    \needspace{5\baselineskip}%
    \nopagebreak[4]%
    \vspace{10pt}
    \noindent\textbf{\textcolor{\fachfarbe}{Aufgabe #1}}
    \nopagebreak[4]%
    \vspace{5pt}
    \nopagebreak[4]%
    \begin{list}{}{\leftmargin=0pt}
    \item[]
}{%
    \end{list}
}

\newcommand{\es}[1]{\textcolor{\fachfarbe}{\textbf{#1}}}

\begin{document}

% ═══════════════════════════════════════════════════════════════
% KOPFBEREICH
% ═══════════════════════════════════════════════════════════════

\begin{center}
    {\Large\bfseries\textcolor{\fachfarbe}{\thema}}\\[0.3em]
    {\large\textcolor{loesung}{\textbf{--- MUSTERLÖSUNG ---}}}
\end{center}

\vspace{5pt}

\noindent
\begin{tabularx}{\textwidth}{|X|X|X|}
\hline
\textbf{Name:} \lsg{Musterschüler/in} &
\textbf{Klasse:} \klasse &
\textbf{Datum:} \datum \\
\hline
\end{tabularx}

\vspace{10pt}

\begin{infobox}
\textbf{Zeit:} \zeit{} | \textbf{Punkte:} \maxpunkte{} BE | \textbf{Hilfsmittel:} keine
\end{infobox}

\vspace{15pt}

% ═══════════════════════════════════════════════════════════════
% AUFGABE 1: Farben zuordnen (10 BE)
% ═══════════════════════════════════════════════════════════════

\begin{aufgabe}{1 -- Farben zuordnen}
Ordne die Farben der richtigen deutschen Übersetzung zu. \punktebox{10}
\end{aufgabe}

\begin{center}
\begin{tabular}{clcl}
\farbbox{rojo} & \lsg{rot} & \farbbox{blanco} & \lsg{weiß} \\[0.8em]
\farbbox{azul} & \lsg{blau} & \farbbox{marron} & \lsg{braun} \\[0.8em]
\farbbox{verde} & \lsg{grün} & \farbbox{gris} & \lsg{grau} \\[0.8em]
\farbbox{amarillo} & \lsg{gelb} & \farbbox{rosa} & \lsg{rosa} \\[0.8em]
\farbbox{negro} & \lsg{schwarz} & \farbbox{naranja} & \lsg{orange} \\
\end{tabular}
\end{center}

\vspace{15pt}

% ═══════════════════════════════════════════════════════════════
% AUFGABE 2: Adjektivangleichung (12 BE)
% ═══════════════════════════════════════════════════════════════

\begin{aufgabe}{2 -- Adjektivangleichung}
Setze das Farbadjektiv in der richtigen Form ein. \punktebox{12}
\end{aufgabe}

\begin{regelbox}
\textbf{Erinnerung:} Adjektive auf -o: roj\textbf{o}/roj\textbf{a}/roj\textbf{os}/roj\textbf{as} \\
Adjektive auf -e/-Konsonant: verde/verdes, azul/azules
\end{regelbox}

\vspace{0.8em}

\noindent
a) La mesa es \lsg{blanca}. \hfill \textit{(weiß)} -- (2P)

\vspace{0.8em}

\noindent
b) El sofá es \lsg{negro}. \hfill \textit{(schwarz)} -- (2P)

\vspace{0.8em}

\noindent
c) Las sillas son \lsg{rojas}. \hfill \textit{(rot)} -- (2P)

\vspace{0.8em}

\noindent
d) Los armarios son \lsg{azules}. \hfill \textit{(blau)} -- (2P)

\vspace{0.8em}

\noindent
e) La cama es \lsg{verde}. \hfill \textit{(grün)} -- (2P)

\vspace{0.8em}

\noindent
f) Las lámparas son \lsg{amarillas}. \hfill \textit{(gelb)} -- (2P)

\vspace{15pt}

% ═══════════════════════════════════════════════════════════════
% AUFGABE 3: Sätze bilden (8 BE)
% ═══════════════════════════════════════════════════════════════

\begin{aufgabe}{3 -- Möbel beschreiben}
Übersetze ins Spanische. Achte auf die Adjektivstellung (Farbe NACH dem Nomen)! \punktebox{8}
\end{aufgabe}

\noindent
a) Der rote Stuhl: \lsg{la silla roja} \hfill (2P)

\vspace{0.8em}

\noindent
b) Das blaue Sofa: \lsg{el sofá azul} \hfill (2P)

\vspace{0.8em}

\noindent
c) Die weißen Tische: \lsg{las mesas blancas} \hfill (2P)

\vspace{0.8em}

\noindent
d) Die schwarzen Schränke: \lsg{los armarios negros} \hfill (2P)

% ═══════════════════════════════════════════════════════════════
% PUNKTEÜBERSICHT
% ═══════════════════════════════════════════════════════════════

\vspace{25pt}
\hrule
\vspace{10pt}

\noindent
\begin{tabularx}{\textwidth}{|X|c|c|c||c|}
\hline
\textbf{Aufgabe} & 1 & 2 & 3 & \textbf{Gesamt} \\
\hline
\textbf{max. Punkte} & 10 & 12 & 8 & \textbf{30} \\
\hline
\textbf{erreicht} & \lsg{10} & \lsg{12} & \lsg{8} & \lsg{30} \\
\hline
\end{tabularx}

\vspace{10pt}

\begin{flushright}
\textbf{Note:} \lsg{1 (14 NP)}
\end{flushright}

\vspace{1em}

\textbf{14-NP-Skala (Sek I):}

\begin{tabular}{|c|c|c|c|c|c|c|c|c|c|c|c|c|c|c|}
\hline
\textbf{NP} & 14 & 13 & 12 & 11 & 10 & 9 & 8 & 7 & 6 & 5 & 4 & 3 & 2 & 1 \\
\hline
\textbf{ab BE} & 30 & 29 & 27 & 26 & 24 & 22 & 20 & 18 & 16 & 14 & 12 & 10 & 8 & 6 \\
\hline
\end{tabular}

\end{document}
